\documentclass[10pt,a4paper]{amsart}
\usepackage[margin=19mm]{geometry}
\usepackage{amsmath,amssymb,mathtools}
\usepackage[scr=rsfs,cal=euler]{mathalfa}
\usepackage{xfrac}
\usepackage[unicode,hidelinks,bookmarks=false]{hyperref}
\newcommand{\ie}{i.e.\ }
\newcommand{\cf}{cf.\ }
\newcommand{\resp}{respectively\ }
\newcommand{\Subsection}{\subsection}
\providecommand{\nobreakdash}{\nobreak\hbox{-}}
\setlength{\emergencystretch}{2em}
\multlinegap0pt
\usepackage{bm}
\usepackage{bbm}
\usepackage{dsfont}
\usepackage{xspace}
\usepackage{cancel}
\usepackage{mathtools}
\usepackage{amsthm}
\makeatletter
\@ifundefined{theorem}{\newtheorem{theorem}{Theorem}}{}
\@ifundefined{lemma}{\newtheorem{lemma}[theorem]{Lemma}}{}
\@ifundefined{proposition}{\newtheorem{proposition}[theorem]{Proposition}}{}
\makeatother
\newcommand{\norm}[1]{\left\lVert #1\right\rVert}
\title[Projection-based approximations for eigenvalue problems of F]{Projection-based approximations for eigenvalue problems of Fredholm integral operators with Green's kernels}
\author{Shashank K. Shukla, Gobinda Rakshit, Akshay S. Rane}
\date{Source: arXiv:2602.16191v1 (CC BY 4.0)}
\begin{document}
\maketitle

\noindent\emph{Source and licence.} Projection-based approximations for eigenvalue problems of Fredholm integral operators with Green's kernels. Version arXiv:2602.16191v1 (CC BY 4.0), distributed under the Creative Commons Attribution 4.0 International licence, as recorded in the arXiv licence metadata for this exact version and shown on the abstract page https://arxiv.org/abs/2602.16191v1. Full rights notice: licenses/arxiv-2602.16191-CC-BY-4.0.txt.\par\smallskip

\noindent\textit{Editorial scope.} Counted unit: the Lemma whose source label is \texttt{lem2} in the article (source0.tex lines 597--614, source label \texttt{lem2}), delivered together with the article's own complete original proof of it (source0.tex lines 615--705, one \texttt{proof} environment of 91 lines). No mathematics is written, repaired or re-derived: statement and proof are the article's own text line for line. Required context retained uncounted: Eq: 002 (lines 89--91); Eq: 005 (lines 247--249); Eq: 007 (lines 313--315); Eq: 008 (lines 317--319); Eq: 009 (lines 321--323); Eq: 0010 (lines 325--327); pro1 (lines 438--448). Every definition, display and label of those lines is retained unchanged. Omitted from this package and not redistributed: 1--88, 92--246, 250--312, 316--316, 320--320, 324--324, 328--437, 449--596, 706--1247. Nothing inside a retained excerpt is omitted or altered: every retained body is delivered exactly as the article's lines are written, so no omission receipt of any kind accompanies this package. Citations: the retained excerpts carry no citation command at all, so no bibliography entry of the article is transcribed and no reference is redistributed. Reference adaptation: none was needed and none was made; every reference inside the delivered unit resolves inside it, so no unresolved reference or citation is produced. Printed slips kept literal: no printed word, symbol, hypothesis or number of the counted statement or of its proof was altered. Layout and class: the article's own class is \texttt{sn-jnl} with packages including graphicx, multirow, amsmath, amssymb, amsfonts, amsthm, mathrsfs, appendix, xcolor, textcomp, manyfoot, booktabs, algorithm, algorithmicx; the wrapper is the lane's pinned \texttt{amsart} editorial wrapper at 10pt on A4, which restates the theorem-like environments the retained text needs and adds the article's own macro definitions that the retained text needs (\texttt{\textbackslash norm}), with extra wrapper packages (bm, bbm, dsfont, xspace, cancel, mathtools). The delivered numbering is the unit's own. Assets: no retained span contains a figure environment, a graphics-inclusion command, a PDF-page inclusion, a table or a TikZ picture; no image file is copied into this package, the package declares no asset dependency and no image file is redistributed with it. Licence: version arXiv:2602.16191v1 of this article is distributed under the Creative Commons Attribution 4.0 International licence, as recorded in the arXiv licence metadata for this exact version and shown on the abstract page https://arxiv.org/abs/2602.16191v1. A full rights notice is delivered at licenses/arxiv-2602.16191-CC-BY-4.0.txt. Characters: every retained excerpt is pure ASCII, so the delivered engine, XeLaTeX with its default Latin Modern text font, reports no missing character.
\begin{equation} \label{Eq: 002}
	(\mathcal{K}x)(s) = \int_{0}^{1} \kappa(s, t)x(t) \, dt, \quad s \in [0, 1], \quad x \in \mathcal{X},
\end{equation}

\begin{equation} \label{Eq: 005}
	\norm{(\mathcal{K}x)^{(2r+1)}}_\infty \leq C_1 \left( \norm{x}_\infty + \norm{x'}_\infty + \cdots + \norm{x^{(2r-1)}}_\infty \right),
\end{equation}

\begin{equation} \label{Eq: 007}
	P_n x = \sum_{j=1}^{n} P_{n,j} x.
\end{equation}

\begin{equation} \label{Eq: 008}
	\norm{P_n} \leq 2 \sum_{\eta=0}^{2r} \norm{L_{\eta}}_{\infty}^2 =C_2. 
\end{equation}

\begin{equation} \label{Eq: 009}
	\norm{(I - P_{n,j}) x}_{\Delta_j, \infty} \leq C_3 \norm{x^{(2r+1)}}_{\Delta_j, \infty} h^{2r+1},
\end{equation}

\begin{equation} \label{Eq: 0010}
	\norm{(I - P_{n})x}_\infty \leq C_3 \norm{x^{(\beta)}}_\infty h^{\beta}.
\end{equation}

\begin{proposition} \label{pro1}
	For $r \geq 1$, if $x \in C^{2r+1}[0,1]$, then
	\begin{eqnarray*}
		\norm{\mathcal{K}(I-P_n)x }_\infty \leq C_5 \norm{x^{(2r+1)}}_{ \infty} h^{2r + 3},
	\end{eqnarray*}
	where $C_5$ is a constant independent of $h$. Also,
	\begin{eqnarray*} 
		\norm{\mathcal{K}(I-P_n)x }_\infty \leq C_6 \norm{x}_\infty h^2,
	\end{eqnarray*}
	for some constant $C_6$ independent of $h$.
\end{proposition}

\begin{lemma} \label{lem2}
	Let $ \mathcal{K}$ be the linear integral operator defined by \eqref{Eq: 002} with a Green's function-type kernel. Let $P_n$ be the orthogonal projection onto the approximating space $\mathcal{X}_n$ defined by \eqref{Eq: 007}. Then
	\begin{eqnarray*}
		\norm{(\mathcal{K} - \mathcal{K}_n^M) \mathcal{K} |_{\mathcal{R}(E)}} =
		\begin{cases}
			O\left(h^{2r+3} \right) & \text{for } r \geq 1, \\
			O\left(h^{3} \right) & \text{for } r = 0,
		\end{cases}
	\end{eqnarray*}
	\begin{eqnarray*}
		\norm{\mathcal{K} (\mathcal{K} - \mathcal{K}_n^M) \mathcal{K} |_{\mathcal{R}(E)}} =
		\begin{cases}
			O\left(h^{2r+5} \right) & \text{for } r \geq 1, \\
			O\left(h^{4} \right) & \text{for } r = 0,
		\end{cases}
	\end{eqnarray*}
	where $\mathcal{K}_n^M= P_n \mathcal{K} + \mathcal{K} P_n - P_n \mathcal{K} P_n.$
\end{lemma}

\begin{proof}
	For any \(\varphi \in \mathcal{R}(E)\) we have  \( \mathcal{K}\varphi \in \mathcal{R}(E)\). Then 
	\begin{align*}
		\norm{(\mathcal{K} - \mathcal{K}_n^M) \mathcal{K} |_{\mathcal{R}(E)}} &= \sup_{\varphi \in \mathcal{R}(E)}\left\{ 	\norm{(\mathcal{K} - \mathcal{K}_n^M) \mathcal{K}\varphi}_\infty : \norm{\varphi}_\infty =1 \right\} \\
		&= \sup_{\varphi \in \mathcal{R}(E)}\left\{ \norm{(I - P_n) \mathcal{K}(I - P_n)\mathcal{K}\varphi}_\infty : \norm{\varphi}_\infty =1 \right\}.
	\end{align*}
	For \( r \geq 1 \), we write  
	\begin{equation*}
		\norm{(I - P_n) \mathcal{K}(I - P_n)\mathcal{K}\varphi}_\infty \leq\norm{I - P_n} \norm{\mathcal{K}(I - P_n)\mathcal{K}\varphi}_\infty.
	\end{equation*}
	From \eqref{Eq: 008}, it follows that \( \norm{I - P_n} \leq 1 +C_2 \). Let \( \varphi \in C^{2r+1}[0,1] \), which implies \( \mathcal{K} \varphi \in C^{2r+1}[0,1] \). Moreover, applying Proposition \ref{pro1}, we have
	\begin{equation*}
		\norm{(I - P_n) \mathcal{K}(I - P_n)\mathcal{K}\varphi}_\infty \leq (1 +C_2)C_5 \norm{(\mathcal{K}\varphi)^{(2r+1)}}_\infty h^{2r+3}.
	\end{equation*}
	Following inequality \eqref{Eq: 005}, we observe that the bound \(\norm{ \left(\mathcal{K}\varphi\right)^{(2r+1)}}_\infty \) does not depend on \(h\). Therefore, 
	\begin{equation} \label{Eq: 0030}
		\norm{(I - P_n) \mathcal{K}(I - P_n)\mathcal{K}\varphi}_\infty = O\left(h^{2r+3}\right).
	\end{equation}
	In the case of \(r=0\), we use \eqref{Eq: 0010} to obtain
	\begin{equation} \label{Eq: 0031}
		\norm{(I - P_n) \mathcal{K}(I - P_n)\mathcal{K}\varphi}_\infty \leq C_3 \norm{\left( \mathcal{K}(I - P_n)y\right)'}_\infty h.
	\end{equation}
	Now, define 
	\begin{equation*}
		\ell_s(t) = \ell(s,t) = \begin{cases}
			D^{(1,0)} \kappa_1(s, t) , & s,t \in \Omega_1, \\
			D^{(1,0)} \kappa_2(s, t) , & s,t \in \Omega_2.
		\end{cases}
	\end{equation*}
	Consider 
	\begin{align*} 
		(\mathcal{K}(I - P_n)y)'(s) &= \int_{0}^{1} \ell_s(t) (I - P_n) y(t) \, dt \notag \\
		&= \sum_{j=1}^{n} \int_{t_{j-1}}^{t_j} \ell_s(t) (I - P_{n,j}) y(t) \, dt \notag \\
		&= \sum_{j=1}^{n} \langle (I - P_{n,j}) \ell_s, (I - P_{n,j}) y \rangle_{\Delta_j}.
	\end{align*}
	Let \( s \in [t_{i-1}, t_i]\), for \(i=1,2,\ldots,n\). Then, we can write
	\begin{equation} \label{Eq: 0032}
		(\mathcal{K}(I - P_n)y)'(s) = \sum_{\substack{j=1 \\  j\ne i} }^{n} \langle (I - P_{n,j}) \ell_s, (I - P_{n,j}) y \rangle_{\Delta_j} + \langle (I - P_{n,i}) \ell_s, (I - P_{n,i}) y \rangle_{\Delta_i}.
	\end{equation}
	For \( j \neq i \), we observe that both \( \ell_s \in C^{2r+1}(\Delta_j) \) and \( x \in C^{2r+1}(\Delta_j)\). Hence, utilizing the estimate \eqref{Eq: 009}, we obtain  
	\begin{equation} \label{Eq: 0033}  
		\bigg|\sum_{\substack{j=1 \\ j\ne i} }^{n} \langle (I - P_{n,j}) \ell_s, (I - P_{n,j}) y \rangle_{\Delta_j} \bigg| \leq (C_3)^2 \norm{\ell_s}_\infty \norm{y}_\infty (n-1) h^{3} \leq  (C_3 C_1)^2 \norm{\varphi}_\infty (n-1) h^{3},
	\end{equation}  
	where \(\norm{\ell_s}_\infty = C_1\) by definition. Also, whenever \(j = i\), \(\ell_s\) is continuous on \( [t_{i-1}, t_i]\), for \(i=1,2,\ldots,n\). Thus, following the proof similar to Proposition \ref{pro1}, we get
	\begin{equation*}
		|\langle (I - P_{n,i}) \ell_s, (I - P_{n,i}) y \rangle_{\Delta_i}| = O\left(h^{2} \right).
	\end{equation*}
	Applying the above estimate along with inequality \eqref{Eq: 0033} in \eqref{Eq: 0032}, we have
	\begin{equation*}
		\norm{	(\mathcal{K}(I - P_n)y)'}_\infty = O\left(h^{2} \right)
	\end{equation*}
	Using above estimate in \eqref{Eq: 0031}, we obtain
	\begin{equation*}
		\norm{(I - P_n) \mathcal{K}(I - P_n)\mathcal{K}\varphi}_\infty = O\left(h^{3} \right).
	\end{equation*}
	The above estimate and \eqref{Eq: 0030} establish the first bound of the lemma.\\
	Now, consider the other bound of lemma as
	\begin{equation*}
		\norm{\mathcal{K}(\mathcal{K} - \mathcal{K}_n^M) \mathcal{K} |_{\mathcal{R}(E)}} 
		= \sup_{\varphi \in \mathcal{R}(E)}\left\{\norm{ \mathcal{K} (I - P_n) \mathcal{K}(I - P_n)\mathcal{K}\varphi}_\infty : \norm{\varphi}_\infty =1 \right\}.
	\end{equation*}
	We can write it as follows
	\begin{equation*}
		\norm{\mathcal{K}(I - P_n) \mathcal{K}(I - P_n)\mathcal{K}\varphi}_\infty \leq\norm{\mathcal{K}(I - P_n)} \norm{\mathcal{K}(I - P_n)\mathcal{K}\varphi}_\infty.
	\end{equation*}
	For any \(x \in C[0,1]\), consider
	\begin{equation*}
		\norm{\mathcal{K}(I - P_n)} = \sup_{x \in C[0,1]}\left\{ \norm{\mathcal{K}(I - P_n)x}_\infty : \norm{x}_\infty =1 \right\}.
	\end{equation*}
	Using Proposition \ref{pro1}, we obtain
	\begin{equation*}
		\norm{\mathcal{K}(I - P_n)} = O\left(h^{2} \right).
	\end{equation*}
	For \(r \geq 1\), let \(\varphi \in C^{2r+1}[0,1]\), it follows that \(\mathcal{K}\varphi \in C^{2r+1}[0,1]\) and using Proposition \ref{pro1}, we have
	\begin{equation*}
		\norm{\mathcal{K}(I - P_n)\mathcal{K}\varphi}_\infty = O\left(h^{2r +3}\right).
	\end{equation*}
	Combining the above two estimates, we obtain 
	\begin{equation}  \label{Eq: 0034}
		\norm{\mathcal{K}(I - P_n) \mathcal{K}(I - P_n)\mathcal{K}\varphi}_\infty = 	O\left(h^{2r+5}\right).
	\end{equation}
	When \(r=0\), using Proposition \ref{pro1}, we obtain
	\begin{equation*}
		\norm{\mathcal{K}(I - P_n)\mathcal{K}\varphi}_\infty = O\left(h^{2}\right).
	\end{equation*}
	Therefore,
	\begin{equation*} 
		\norm{\mathcal{K}(I - P_n) \mathcal{K}(I - P_n)\mathcal{K}\varphi}_\infty = 	O\left(h^{4}\right).
	\end{equation*}
	The above estimate along with \eqref{Eq: 0034} establish the desired second bound of the lemma, thereby completing the proof.
\end{proof}

\end{document}
